\documentclass[10pt]{article}

% Pacotes extras necessários
\usepackage{amsmath}
\usepackage[lmargin=0.5in, rmargin=0.5in, tmargin=0.5in, bmargin=0.5in, includehead, includefoot]{geometry}
\usepackage{amsfonts}
\usepackage[utf8]{inputenc}
\usepackage[portuguese]{babel}
\usepackage{graphicx}
\usepackage{fancyhdr}
\usepackage{setspace}
\usepackage{listings}
\usepackage{url}
\usepackage{enumitem}
\usepackage{appendix}
\usepackage{subcaption}
\usepackage[dvipsnames]{xcolor}
\usepackage{tikz}
\usepackage{makecell}
\usetikzlibrary{positioning}

% Define estilo para listagens de código
\lstdefinestyle{mystyle}{
    language=Matlab,
    inputencoding=latin1,
    extendedchars=true,
    backgroundcolor=\color{white},
    basicstyle=\ttfamily\footnotesize,
    numbers=left,
    numberstyle=\tiny\color{gray},
    numbersep=5pt,
    tabsize=2,
    breaklines=true,
    keywordstyle=\color{blue},
    commentstyle=\color{green!60!black},
    stringstyle=\color{orange},
    frame=single,
    keepspaces=true,
    showspaces=false,
    showstringspaces=false,
    showtabs=false,
}
\lstset{style=mystyle}

\graphicspath{{./images/}}

\setlength{\parindent}{0pt}
\setstretch{1.5}

\DeclareMathOperator{\sen}{sen}
\DeclareMathOperator{\sinc}{sinc}
\newcommand{\Lap}[1]{\mathcal{L}\left\{#1\right\}}
\newcommand{\bm}[1]{\boldsymbol{#1}}

% Cabeçalho e rodapé em todas as páginas
\pagestyle{fancy}
\fancyhf{}
\fancyhead[L]{Relatório das simulações}
\fancyfoot[C]{\thepage}

% Dados do Grupo
\title{
    Trabalho Nº 07 -  Combined M-MRAC + LS \\
    \large COE603 - Controle Adaptativo
}
\author{
    Caio Cesar Leal Verissimo - 119046624 \\
    Leonardo Soares da Costa Tanaka - 121067652 \\
    Lincoln Rodrigues Proença - 121076407 \\
    Engenharia de Controle e Automação - UFRJ \\
    Rio de Janeiro, Rio de Janeiro, Brasil \\
    Julho de 2025
}
\date{}

\begin{document}

\maketitle

% Geração automática do sumário
\tableofcontents
\newpage

\section{Descrição do Algoritmo}

\subsection{Introdução ao Algoritmo}

\subsection{Análise de Estabilidade}

\subsection{Resumo do Algoritmo}

\begin{table}[htbp]
\centering
\begin{tabular}{|>{\bfseries}m{4cm}|m{7cm}|c|}
\hline
\textcolor{blue}{Subsistema} & \textcolor{blue}{Equação} & \textcolor{blue}{Ordem} \\
\hline
Planta & $y = P(s)\,u$ & $n$ \\
\hline
Modelo & $y_m = M(s)\,r$ & $n$ \\
\hline
Erro de rastreamento & $e_a = \text{sign}(k_p)(y - y_m)$ & \\
\hline
Filtros SV & 
$\begin{aligned}
\dot{\omega}_1 &= A_f \omega_1 + b_f u \\
\dot{\omega}_2 &= A_f \omega_2 + b_f y
\end{aligned}$ & $n-1$ \\
\hline
Regressor & $\omega^T = \begin{bmatrix} \omega_1^T & y & \omega_2^T & r \end{bmatrix}$ & \\
\hline
Filtro $\xi$ & 
$\begin{aligned}
\xi &= -\ell_0 \xi + \omega, \quad \ell_0 > a_m
\end{aligned}$ & $2n$ \\
\hline
Lei de controle & $u = \theta^T \omega + \dot{\theta}^T \xi$ & \\
\hline
Lei de adaptação & 
$\begin{aligned}
\dot{\theta} &= -\Gamma \xi e_a - \sigma \Delta, \quad \Gamma = \Gamma^T > 0, \quad \sigma > 0 \\
\Delta &= \theta - \psi
\end{aligned}$ & $2n$ \\
\hline
Filtros auxiliares & 
$\begin{aligned}
\dot{\zeta} &= -\ell_0 \zeta + u \\
\dot{\varphi} &= -\ell_0 \varphi + e_0
\end{aligned}$ & $1$ \\
\hline
Predição & 
$\begin{aligned}
\hat{\zeta} &= \psi^T \phi \\
\phi &= \xi + \left[ \mathbf{0}^T \quad (e_0 - \alpha \varphi) \right]^T, \quad \mathbf{0} \in \mathbb{R}^{2n-1}
\end{aligned}$ & \\
\hline
Erro de predição & $\varepsilon = \hat{\zeta} - \zeta$ & \\
\hline
Estimador LS & 
$\begin{aligned}
\dot{\psi} &= -R \left( \tau \phi \varepsilon - \frac{\sigma}{\beta} \Gamma^{-1} \Delta \right), \quad \tau > \frac{1}{2}, \quad \beta > 0 \\
\dot{R} &= - \frac{R \phi \phi^T R}{m^2}, \quad R(0) = R^T(0) > 0 \\
m^2 &= 1 + \kappa \phi^T R \phi, \quad \kappa \geq 0
\end{aligned}$ & 
$\begin{aligned}
2n \\
\textcolor{red}{4n^2}
\end{aligned}$ \\
\hline
\end{tabular}
\caption{Resumo dos subsistemas, equações e ordens no controle adaptativo}
\vspace{1em}
\textbf{Dimensão total do sistema:} \quad
\fcolorbox{magenta}{white}{$\displaystyle N = 10n + \textcolor{red}{4n^2}$}
\end{table}

\section{Resultados das Simulações}

\section{Discussão dos Resultados}


\end{document}
